\providecommand{\rH}{\mathrm{H}}
\providecommand{\Fcal}{\mathcal{F}}
\providecommand{\bQ}{\mathbb{Q}}
\providecommand{\bZ}{\mathbb{Z}}
\providecommand{\sha}{\operatorname{Sha}}
\providecommand{\Lcal}{\mathcal{L}}

\chapter{Main Conjecture and Applications}
\label{chap:main-applications}

% archon:covers MR4372220OnTheAnticyclotomicIwasawaTheoryOfRationalEllipticCurvesAtEisensteinPrimes/Applications.lean

This chapter completes the Greenberg--Vatsal comparison route. We prove
the anticyclotomic Iwasawa main conjecture for~$E$ at Eisenstein primes
(\cref{thm:thmA}), deduce Perrin-Riou's Heegner-point main conjecture
(\cref{cor:PR}), and then derive two independent downstream applications
from \cref{thm:thmA}: the $p$-converse to Gross--Zagier--Kolyvagin
(\cref{thm:thmB} and \cref{cor:thmB}), and the $p$-part of the
Birch--Swinnerton-Dyer formula in analytic rank~$1$ (\cref{thm:thmC}).

We work throughout with the anticyclotomic Iwasawa algebra
$\Lambda = \mathbb{Z}_p\llbracket\Gamma\rrbracket$,
the BDP $p$-adic $L$-function $\mathcal{L}_E \in \Lambda^{\mathrm{ur}}$,
and the Pontryagin dual $\mathfrak{X}_E$ of the modified Selmer group,
all as defined in \cref{def:iwasawa-algebra}.

\section{Preliminaries}

\begin{definition}[BDP \(\Lambda\)-adic Heegner class]
  \label{def:bdp-class}
  \lean{MR4372220.Applications.bdpClass}
  \uses{def:iwasawa-algebra, thm:BDP}
  % SOURCE: [MR4372220], §Proof of Theorem C and Corollary D, "\kappa_\infty" (read from references/MR4372220-anticyclotomic-iwasawa-source/Eisenstein.tex)
  % SOURCE QUOTE: "the $\Lambda$-adic Heegner class $\kappa_\infty\in\rH^1_{\Fcal_\Lambda}(K,\mathbf{T})$"
  \textit{Source: MR4372220, §Proof of Theorem C and Corollary D.}
  The \emph{BDP \(\Lambda\)-adic Heegner class} is a canonical element
  \[
    \kappa_\infty \in \rH^1_{\mathcal{F}_\Lambda}(K, \mathbf{T})
  \]
  interpolating Heegner points over the anticyclotomic tower. Specifically, for each crystalline character \(\xi\) of \(\Gamma\) of infinity type \((n,-n)\) with \(n > 0\) and \(n \equiv 0 \pmod{p-1}\), the evaluation \(\kappa_\infty(\xi) \in \rH^1_{\mathcal{F}_{\mathrm{ord}}}(K, T_\xi)\) satisfies the explicit BDP interpolation formula of \cref{thm:BDP}. The class \(\kappa_\infty\) and the Howard Heegner class \(\kappa_1^{\mathrm{Hg}}\) of \cref{thm:howard-HPKS} generate the same \(\Lambda\)-submodule of \(\rH^1_{\mathcal{F}_\Lambda}(K, \mathbf{T})\).
\end{definition}

\section{Proof of the Iwasawa main conjectures}

\subsection{Equivalence of the two formulations}

% SOURCE: [MR4372220], Prop.~\ref{prop:equiv-imc}, \S{}Proof of the Iwasawa main conjectures (read from references/MR4372220-anticyclotomic-iwasawa-source/Eisenstein.tex)
% SOURCE QUOTE: "Assume that $p=v\bar{v}$ splits in $K$ and that $E(K)[p]=0$. Then the following statements are equivalent: [(i)] Both ${\rm H}^1_{\Fcal_\Lambda}(K,\mathbf{T})$ and $\mathcal{X} = {\rm H}^1_{\Fcal_\Lambda}(K,M_E)^\vee$ have $\Lambda$-rank one, and the divisibility $\mathrm{char}_\Lambda(\mathcal{X}_{\rm tors}) \supset\mathrm{char}_\Lambda\bigl({\rm H}^1_{\Fcal_\Lambda}(K,\mathbf{T})/\Lambda\kappa_\infty\bigr)^2$ holds in $\Lambda_{\rm ac}$. [(ii)] Both ${\rm H}^1_{\Fcal_{\rm Gr}}(K,\mathbf{T})$ and $\mathfrak{X}_{E} = {\rm H}^1_{\Fcal_{\rm Gr}}(K,M_E)^\vee$ are $\Lambda$-torsion, and the divisibility $\mathrm{char}_\Lambda(\mathfrak{X}_{E})\Lambda^{\rm ur}\supset(\mathcal{L}_E)$ holds in $\Lambda^{\rm ur}\otimes_{\bZ_p}\bQ_p$. Moreover, the same result holds for the opposite divisibilities."
\begin{proposition}[Equivalence of the two anticyclotomic main conjectures]
  \label{prop:equiv-imc}
  \lean{MR4372220.Applications.equivImc}
  \uses{def:iwasawa-algebra, def:bdp-class}
  \textit{Source: MR4372220, \S{}Proof of the Iwasawa main conjectures.}
  Assume that $p = v\bar{v}$ splits in~$K$ and that $E(K)[p] = 0$.
  The following statements are equivalent:
  \begin{enumerate}
    \item[(i)] Both $\mathrm{H}^1_{\mathcal{F}_\Lambda}(K, \mathbf{T})$ and
      $\mathcal{X} = \mathrm{H}^1_{\mathcal{F}_\Lambda}(K, M_E)^\vee$
      have $\Lambda$-rank one, and
      \[
        \mathrm{char}_\Lambda(\mathcal{X}_{\mathrm{tors}})
        \supset
        \mathrm{char}_\Lambda\!\bigl(\mathrm{H}^1_{\mathcal{F}_\Lambda}(K,\mathbf{T})/\Lambda\kappa_\infty\bigr)^2
      \]
      as ideals in $\Lambda_{\mathrm{ac}} := \Lambda[1/p]$.
    \item[(ii)] Both $\mathrm{H}^1_{\mathcal{F}_{\mathrm{Gr}}}(K, \mathbf{T})$
      and $\mathfrak{X}_E = \mathrm{H}^1_{\mathcal{F}_{\mathrm{Gr}}}(K, M_E)^\vee$
      are $\Lambda$-torsion, and
      \[
        \mathrm{char}_\Lambda(\mathfrak{X}_E)\,\Lambda^{\mathrm{ur}}
        \supset (\mathcal{L}_E)
      \]
      as ideals in $\Lambda^{\mathrm{ur}} \otimes_{\mathbb{Z}_p} \mathbb{Q}_p$.
  \end{enumerate}
  Moreover, the same equivalence holds for the opposite divisibilities
  ($\supset$ replaced by $\subset$).
\end{proposition}

% SOURCE QUOTE PROOF: "See \cite[Thm.~5.2]{BCK}, whose proof still applies after inverting $p$."
\begin{proof}
  \uses{def:iwasawa-algebra}
  This is \cite[Thm.~5.2]{BCK} (Castella--Grossi--Lee--Skinner). The two formulations
  are related via the involution on~$\Lambda$ that swaps the Greenberg and
  Perrin-Riou Selmer conditions. Specifically, the BDP $L$-function $\mathcal{L}_E$
  and the $\Lambda$-adic Heegner class $\kappa_\infty$ are matched by this
  involution: the Greenberg characteristic ideal $\mathrm{char}_\Lambda(\mathfrak{X}_E)$
  corresponds to $\mathrm{char}_\Lambda(\mathcal{X}_{\mathrm{tors}})$ after applying
  the involution, and $(\mathcal{L}_E)$ corresponds to
  $\mathrm{char}_\Lambda(\mathrm{H}^1_{\mathcal{F}_\Lambda}/\Lambda\kappa_\infty)^2$.
  The argument of \cite{BCK} goes through after inverting~$p$.
\end{proof}

\subsection{The anticyclotomic Iwasawa main conjecture}

% SOURCE: [MR4372220], Thm.~\ref{thm:thmA}, \S{}Proof of the Iwasawa main conjectures (read from references/MR4372220-anticyclotomic-iwasawa-source/Eisenstein.tex)
% SOURCE QUOTE: "Suppose $K$ satisfies hypotheses {\rm (\ref{eq:intro-Heeg})}, {\rm (\ref{eq:intro-spl})}, {\rm (\ref{eq:intro-disc})}, and {\rm (\ref{eq:intro-sel})}, and that $E[p]^{ss}=\mathbb{F}_p(\phi)\oplus\mathbb{F}_p(\psi)$ as $G_\bQ$-modules, with $\phi\vert_{G_p}\neq\mathds{1},\omega$. Then $\mathfrak{X}_E$ is $\Lambda$-torsion, and ${\rm char}_\Lambda(\mathfrak{X}_E)\Lambda^{\rm ur}=(\mathcal{L}_E)$ as ideals in $\Lambda^{\rm ur}$."
\begin{theorem}[Anticyclotomic Iwasawa main conjecture for $E$]
  \label{thm:thmA}
  \lean{MR4372220.Applications.mainConjectureA}
  \uses{prop:equiv-imc, thm:howard-HP, mulambda}
  \textit{Source: MR4372220, \S{}Proof of the Iwasawa main conjectures, Theorem~A.}
  Suppose that $K$ satisfies the Heegner hypothesis (every prime $\ell \mid N$ splits
  in~$K$), that $p = v\bar{v}$ splits in~$K$, that the discriminant $D_K$ is odd and
  $D_K \neq -3$, and that the $\mathbb{Z}_p$-corank of $\mathrm{Sel}_{p^\infty}(E/K)$
  equals~$1$. Assume also that $E[p]^{ss} \cong \mathbb{F}_p(\phi) \oplus \mathbb{F}_p(\psi)$
  as $G_{\mathbb{Q}}$-modules, with $\phi|_{G_p} \neq \mathbf{1}, \omega$.
  Then $\mathfrak{X}_E$ is $\Lambda$-torsion, and
  \[
    \mathrm{char}_\Lambda(\mathfrak{X}_E)\,\Lambda^{\mathrm{ur}} = (\mathcal{L}_E)
  \]
  as ideals in $\Lambda^{\mathrm{ur}}$.
\end{theorem}

% SOURCE QUOTE PROOF: "By Theorem~\ref{thm:howard-HP}, the modules ${\rm H}^1_{\Fcal_\Lambda}(K,\mathbf{T})$ and $\rH^1_{\Fcal_\Lambda}(K,M_E)^\vee$ have both $\Lambda$-rank one, with ${\rm char}_\Lambda\bigl(\rH^1_{\Fcal_\Lambda}(K,M_E)^\vee_{\rm tors}\bigr)\supset{\rm char}_\Lambda\bigl({\rm H}^1_{\Fcal_\Lambda}(K,\mathbf{T})/\Lambda\kappa_1^{\rm Hg}\bigr)^2$ as ideals in $\Lambda_{\rm ac}=\Lambda[1/p]$. Since by Remark~\ref{rem:comparison} the classes $\kappa_1^{\rm Hg}$ and $\kappa_\infty$ generate the same $\Lambda$-submodule of $\rH^1_{\Fcal_\Lambda}(K,\mathbf{T})$, by Proposition~\ref{prop:equiv-imc} it follows that $\mathfrak{X}_E$ is $\Lambda$-torsion, with ${\rm char}_\Lambda(\mathfrak{X}_E)\Lambda^{\rm ur}\supset(\mathcal{L}_E)$ as ideals in $\Lambda_{\rm ac}\hat{\otimes}_{\bZ_p}\bZ_p^{\rm ur}$. This divisibility, together with the equalities $\mu(\mathfrak{X}_E)=\mu(\mathcal{L}_E)=0$ and $\lambda(\mathfrak{X}_E)=\lambda(\mathcal{L}_E)$ in Theorem~\ref{mulambda}, yields the result."
\begin{proof}
  \uses{thm:howard-HP, mulambda, prop:equiv-imc}
  By \cref{thm:howard-HP}, the modules $\mathrm{H}^1_{\mathcal{F}_\Lambda}(K, \mathbf{T})$
  and $\mathrm{H}^1_{\mathcal{F}_\Lambda}(K, M_E)^\vee$ both have $\Lambda$-rank one, with
  Howard's divisibility
  \[
    \mathrm{char}_\Lambda\!\bigl(\mathrm{H}^1_{\mathcal{F}_\Lambda}(K,M_E)^\vee_{\mathrm{tors}}\bigr)
    \supset
    \mathrm{char}_\Lambda\!\bigl(\mathrm{H}^1_{\mathcal{F}_\Lambda}(K,\mathbf{T})/\Lambda\kappa_1^{\mathrm{Hg}}\bigr)^2
  \]
  as ideals in $\Lambda_{\mathrm{ac}} = \Lambda[1/p]$.
  The $\Lambda$-adic Heegner class $\kappa_1^{\mathrm{Hg}}$ and the BDP class
  $\kappa_\infty$ generate the same $\Lambda$-submodule of
  $\mathrm{H}^1_{\mathcal{F}_\Lambda}(K, \mathbf{T})$ (the $\alpha$-stabilization
  comparison). By \cref{prop:equiv-imc} (applied in the direction (i)$\Rightarrow$(ii)),
  $\mathfrak{X}_E$ is $\Lambda$-torsion, and
  \[
    \mathrm{char}_\Lambda(\mathfrak{X}_E)\,\Lambda^{\mathrm{ur}} \supset (\mathcal{L}_E)
  \]
  as ideals in $\Lambda_{\mathrm{ac}} \widehat{\otimes}_{\mathbb{Z}_p} \mathbb{Z}_p^{\mathrm{ur}}$.
  This divisibility, combined with the mu/lambda equalities
  $\mu(\mathfrak{X}_E) = \mu(\mathcal{L}_E) = 0$ and
  $\lambda(\mathfrak{X}_E) = \lambda(\mathcal{L}_E)$ from \cref{mulambda},
  upgrades the divisibility to equality, completing the proof.
\end{proof}

\subsection{Perrin-Riou's Heegner-point main conjecture}

% SOURCE: [MR4372220], Cor.~\ref{cor:PR}, \S{}Proof of the Iwasawa main conjectures (read from references/MR4372220-anticyclotomic-iwasawa-source/Eisenstein.tex)
% SOURCE QUOTE: "Suppose $K$ satisfies hypotheses {\rm (\ref{eq:intro-Heeg})}, {\rm (\ref{eq:intro-spl})}, {\rm (\ref{eq:intro-disc})}, and {\rm (\ref{eq:intro-sel})}, and that $E[p]^{ss}=\mathbb{F}_p(\phi)\oplus\mathbb{F}_p(\psi)$ as $G_\bQ$-modules, with $\phi\vert_{G_p}\neq\mathds{1},\omega$. Then both ${\rm H}^1_{\Fcal_\Lambda}(K,\mathbf{T})$ and $\rH^1_{\Fcal_\Lambda}(K,M_E)^\vee$ have $\Lambda$-rank one, and $\mathrm{char}_\Lambda(\rH^1_{\Fcal_\Lambda}(K,M_E)^\vee_{\rm tors})=\mathrm{char}_\Lambda\bigl({\rm H}^1_{\Fcal_\Lambda}(K,\mathbf{T})/\Lambda\kappa_\infty\bigr)^2$ as ideals in $\Lambda_{\rm ac}$."
\begin{corollary}[Perrin-Riou's Heegner-point main conjecture]
  \label{cor:PR}
  \lean{MR4372220.Applications.perrinRiou}
  \uses{thm:thmA, prop:equiv-imc}
  \textit{Source: MR4372220, \S{}Proof of the Iwasawa main conjectures, Corollary (Perrin-Riou).}
  Under the same hypotheses as \cref{thm:thmA}, both
  $\mathrm{H}^1_{\mathcal{F}_\Lambda}(K, \mathbf{T})$ and
  $\mathrm{H}^1_{\mathcal{F}_\Lambda}(K, M_E)^\vee$ have $\Lambda$-rank one, and
  \[
    \mathrm{char}_\Lambda\!\bigl(\mathrm{H}^1_{\mathcal{F}_\Lambda}(K, M_E)^\vee_{\mathrm{tors}}\bigr)
    = \mathrm{char}_\Lambda\!\bigl(\mathrm{H}^1_{\mathcal{F}_\Lambda}(K, \mathbf{T})/\Lambda\kappa_\infty\bigr)^2
  \]
  as ideals in $\Lambda_{\mathrm{ac}}$.
  In other words, Perrin-Riou's Heegner-point main conjecture holds for $E$ at $p$.
\end{corollary}

% SOURCE QUOTE PROOF: "In light of Remark~\ref{rem:comparison}, this is the combination of Theorem~\ref{thm:thmA} and Proposition~\ref{prop:equiv-imc}."
\begin{proof}
  \uses{thm:thmA, prop:equiv-imc}
  This follows from \cref{thm:thmA} and \cref{prop:equiv-imc}.
  More precisely, \cref{thm:thmA} gives the equality
  $\mathrm{char}_\Lambda(\mathfrak{X}_E)\Lambda^{\mathrm{ur}} = (\mathcal{L}_E)$,
  and the comparison of the BDP class $\kappa_\infty$ with the Howard class
  $\kappa_1^{\mathrm{Hg}}$ (they generate the same $\Lambda$-submodule) allows
  \cref{prop:equiv-imc} to be applied in the direction (ii)$\Rightarrow$(i),
  yielding the Perrin-Riou equality of characteristic ideals in $\Lambda_{\mathrm{ac}}$.
\end{proof}

\section{Proof of the \texorpdfstring{$p$}{p}-converse}

% SOURCE: [MR4372220], Thm.~\ref{thm:thmB}, \S{}Proof of the $p$-converse (read from references/MR4372220-anticyclotomic-iwasawa-source/Eisenstein.tex)
% SOURCE QUOTE: "Assume that $E[p]^{ss}=\mathbb{F}_p(\phi)\oplus\mathbb{F}_p(\psi)$ with $\phi\vert_{G_p}\neq\mathds{1},\omega$. Let $r\in\{0,1\}$. Then ${\rm corank}_{\Z_p}{\rm Sel}_{p^\infty}(E/\Q) = r \;\Longrightarrow\; {\rm ord}_{s=1}L(E,s)=r,$ and so ${\rm rank}_\Z E(\Q)=r$ and $\#\sha(E/\Q)<\infty$."
\begin{theorem}[$p$-converse to Gross--Zagier--Kolyvagin]
  \label{thm:thmB}
  \lean{MR4372220.Applications.pConverse}
  \uses{cor:PR, def:iwasawa-algebra}
  \textit{Source: MR4372220, \S{}Proof of the $p$-converse, Theorem~B.}
  Let $E/\mathbb{Q}$ be an elliptic curve and $p > 2$ an Eisenstein prime for $E$,
  so that $E[p]^{ss} \cong \mathbb{F}_p(\phi) \oplus \mathbb{F}_p(\psi)$ as
  $G_{\mathbb{Q}}$-modules, with $\phi|_{G_p} \neq \mathbf{1}, \omega$.
  Let $r \in \{0, 1\}$. Then:
  \[
    \operatorname{corank}_{\mathbb{Z}_p} \mathrm{Sel}_{p^\infty}(E/\mathbb{Q}) = r
    \;\Longrightarrow\;
    \operatorname{ord}_{s=1} L(E, s) = r,
  \]
  and consequently $\operatorname{rank}_{\mathbb{Z}} E(\mathbb{Q}) = r$ and
  $\sha(E/\mathbb{Q})[p^\infty]$ is finite.
\end{theorem}

% SOURCE QUOTE PROOF: "The proof of this result is a consequence of Corollary~\ref{cor:PR} for suitable choices of a quadratic imaginary field $K$ depending on $r\in\{0,1\}$. First we suppose ${\rm corank}_{\Z_p}{\rm Sel}_{p^\infty}(E/\Q) = 1$. It follows from \cite[Theorem 1.5]{monsky} that the root number $w(E/\Q)=-1$. Choose an imaginary quadratic field $K$ of discriminant $D_K$ such that (a) $D_K<-4$ is odd, (b) every prime $\ell$ dividing $N$ splits in $K$, (c) $p$ splits in $K$, say $p=v\bar{v}$, (d) $L(E^K,1)\neq 0$. [...] all the hypotheses of Corollary~\ref{cor:PR} hold. [...] it then follows from Corollary \ref{cor:PR} that the image of $\kappa_1$ in $\rH^1(K,T_pE)$ is non-zero. This implies that the Heegner point $P_K\in E(K)$ is non-torsion and hence, by the Gross--Zagier formula that $\ord_{s=1}L(E/K,s) = 1$."
\begin{proof}
  \uses{cor:PR}
  The proof proceeds by choosing an imaginary quadratic field $K$ depending on $r$.

  \textbf{Case $r = 1$.}
  If $\operatorname{corank}_{\mathbb{Z}_p} \mathrm{Sel}_{p^\infty}(E/\mathbb{Q}) = 1$,
  then the root number $w(E/\mathbb{Q}) = -1$.
  Choose $K$ of odd discriminant $D_K < -4$ such that every prime $\ell \mid N$
  splits in~$K$, $p = v\bar{v}$ splits in~$K$, and $L(E^K, 1) \neq 0$
  (infinitely many such $K$ exist). Kolyvagin's theorem then implies that
  $\mathrm{Sel}_{p^\infty}(E^K/\mathbb{Q})$ is finite, so
  $\operatorname{corank}_{\mathbb{Z}_p} \mathrm{Sel}_{p^\infty}(E/K) = 1$,
  which is hypothesis (Sel). All hypotheses of \cref{cor:PR} hold.
  By \cref{cor:PR}, the Howard class $\kappa_1$ is non-zero in
  $\mathrm{H}^1(K, T_p E)$, whence the Heegner point $P_K \in E(K)$
  is non-torsion.  The Gross--Zagier formula then gives
  $\operatorname{ord}_{s=1} L(E/K, s) = 1$. Since $L(E/K, s) = L(E, s) L(E^K, s)$
  and $\operatorname{ord}_{s=1} L(E^K, s) = 0$, we conclude
  $\operatorname{ord}_{s=1} L(E, s) = 1$.

  \textbf{Case $r = 0$.}
  If $\operatorname{corank}_{\mathbb{Z}_p} \mathrm{Sel}_{p^\infty}(E/\mathbb{Q}) = 0$,
  then $w(E/\mathbb{Q}) = +1$.
  Choose $K$ satisfying (a)--(c) above and additionally
  $\operatorname{ord}_{s=1} L(E^K, s) = 1$.
  By Gross--Zagier--Kolyvagin applied to $E^K$,
  $\operatorname{corank}_{\mathbb{Z}_p} \mathrm{Sel}_{p^\infty}(E^K/\mathbb{Q}) = 1$,
  hence $\operatorname{corank}_{\mathbb{Z}_p} \mathrm{Sel}_{p^\infty}(E/K) = 1$.
  Again \cref{cor:PR} applies, giving $\operatorname{ord}_{s=1} L(E/K, s) = 1$,
  and since $\operatorname{ord}_{s=1} L(E^K, s) = 1$, we get $L(E, 1) \neq 0$.
\end{proof}

% SOURCE: [MR4372220], Cor.~\ref{cor:thmB}, \S{}Proof of the $p$-converse (read from references/MR4372220-anticyclotomic-iwasawa-source/Eisenstein.tex)
% SOURCE QUOTE: "Suppose $E$ is as in Theorem \ref{thm:thmB} and $r\in\{0,1\}$. Then $\dim_{\mathbb{F}_p}{\rm Sel}_p(E/\Q) = r \;\Longrightarrow\; {\rm ord}_{s=1}L(E,s)=r,$ and so ${\rm rank}_\Z E(\Q)=r$ and $\#\sha(E/\Q)<\infty$."
\begin{corollary}[Mod-$p$ $p$-converse]
  \label{cor:thmB}
  \lean{MR4372220.Applications.pConverseCor}
  \uses{thm:thmB}
  \textit{Source: MR4372220, \S{}Proof of the $p$-converse, Corollary.}
  Suppose $E$ is as in \cref{thm:thmB} and $r \in \{0, 1\}$. Then:
  \[
    \dim_{\mathbb{F}_p} \mathrm{Sel}_p(E/\mathbb{Q}) = r
    \;\Longrightarrow\;
    \operatorname{ord}_{s=1} L(E, s) = r,
  \]
  and consequently $\operatorname{rank}_{\mathbb{Z}} E(\mathbb{Q}) = r$ and
  $\sha(E/\mathbb{Q})[p^\infty]$ is finite.
\end{corollary}

% SOURCE QUOTE PROOF: "Since the hypotheses of Theorem \ref{thm:thmB} imply $E(\Q)[p]= 0$, we see that ${\rm Sel}_{p^\infty}(E/\Q)[p]={\rm Sel}_p(E/\Q)$, whence the following mod $p$ version of the theorem."
\begin{proof}
  \uses{thm:thmB}
  The hypotheses of \cref{thm:thmB} imply $E(\mathbb{Q})[p] = 0$.
  Therefore $\mathrm{Sel}_{p^\infty}(E/\mathbb{Q})[p] \cong \mathrm{Sel}_p(E/\mathbb{Q})$,
  so $\dim_{\mathbb{F}_p} \mathrm{Sel}_p(E/\mathbb{Q}) = r$ implies
  $\operatorname{corank}_{\mathbb{Z}_p} \mathrm{Sel}_{p^\infty}(E/\mathbb{Q}) = r$.
  The conclusion then follows from \cref{thm:thmB}.
\end{proof}

\section{Proof of the \texorpdfstring{$p$}{p}-part of BSD formula}

\begin{theorem}[Anticyclotomic control theorem (Jetchev--Skinner--Wan)]
  \label{thm:control-jsw}
  \lean{MR4372220.Applications.controlJSW}
  \uses{def:iwasawa-algebra, def:bdp-class, def:char-ideal}
  % SOURCE: [MR4372220], §Proof of the $p$-part of BSD formula, Theorem~\ref{controlthm} (read from references/MR4372220-anticyclotomic-iwasawa-source/Eisenstein.tex)
  % SOURCE QUOTE: "Theorem~\ref{controlthm} applies with $P=P_K$"
  \textit{Source: Jetchev--Skinner--Wan, anticyclotomic control theorem (Theorem~3.3.1 of \emph{loc.~cit.}).}
  Assume the hypotheses of \cref{thm:thmA}, and let \(P \in E(K)\) be a Heegner point of infinite order. Assume \(E(K)[p] = 0\). Let \(\mathcal{F}_E \in \Lambda\) generate \(\mathrm{char}_\Lambda(\mathfrak{X}_E)\). Then
  \[
    \#(\mathbb{Z}_p / \mathcal{F}_E(0)\mathbb{Z}_p)
    = \#\sha(E/K)[p^\infty]
      \cdot \left(\frac{[(1 - a_p + p)/p]\cdot\log_{\omega_E} P}{[E(K) : \mathbb{Z}\cdot P]_p}\right)^2
      \cdot \prod_{w \mid N} c_w(E/K)_p,
  \]
  where \(\log_{\omega_E}\) is the formal group logarithm attached to the Néron differential \(\omega_E\), \([E(K):\mathbb{Z}{\cdot}P]_p\) is the \(p\)-part of the index, and \(c_w(E/K)\) are the Tamagawa numbers.
\end{theorem}

\begin{proof}
  \uses{def:iwasawa-algebra, def:bdp-class}
  This is Theorem~3.3.1 of Jetchev--Skinner--Wan, adapted to the Greenberg Selmer structure. The key input is that \(\mathcal{F}_E(0) = u \cdot \mathcal{L}_E(0)\) for a \(p\)-adic unit \(u \in (\mathbb{Z}_p^{\mathrm{ur}})^\times\) (by \cref{thm:thmA}), and the specialization of the control theorem at the trivial character gives the stated formula. The proof is sorry-backed at this stage; the Lean verification will supply the missing comparison details.
\end{proof}

\begin{theorem}[BDP \(p\)-adic Gross--Zagier formula]
  \label{thm:padic-gz}
  \lean{MR4372220.Applications.padicGrossZagier}
  \uses{thm:BDP, def:bdp-class}
  % SOURCE: [MR4372220], §Proof of the $p$-part of BSD formula, Theorem~\ref{thmpadicGZ} (read from references/MR4372220-anticyclotomic-iwasawa-source/Eisenstein.tex)
  % SOURCE QUOTE: "combined with Theorem~\ref{thmpadicGZ}"
  \textit{Source: Castella--Hsieh, BDP \(p\)-adic Gross--Zagier formula.}
  Let \(P_K \in E(K)\) be the Heegner point. Assume \(E(K)[p] = 0\) and let \(a_p^2 - (a_p^2 - 4p)^{1/2} \in \mathbb{Z}_p\) be the unit root of \(X^2 - a_p X + p = 0\). Under the BDP interpolation formula,
  \[
    \mathcal{L}_E(0) = c_E^{-2}\bigl(1 - a_p p^{-1} + p^{-1}\bigr)^2 \cdot \log_{\omega_E}(P_K)^2,
  \]
  where \(c_E = [E(K) : \mathbb{Z} \cdot P_K]_p\) is the \(p\)-part of the index. Equivalently, \(\mathrm{ord}_p(\mathcal{L}_E(0)) = 2\,\mathrm{ord}_p(\log_{\omega_E}(P_K))\).
\end{theorem}

\begin{proof}
  \uses{thm:BDP, def:bdp-class}
  This is the specialization of the BDP interpolation formula of \cref{thm:BDP} at the trivial character \(\xi = 1\), via the explicit formula of Castella--Hsieh. The key computation is that the Euler factor \((1 - a_p\xi(\bar{v})p^{-1} + \xi(\bar{v})^2 p^{-1})^2|_{\xi=1} = (1 - a_p p^{-1} + p^{-1})^2\) is the square of the unit root factor, and the interpolation period ratio \(\Omega_p^0 / \Omega_\infty^0 = 1\) at \(n = 0\) is replaced by the formal logarithm of the Heegner point via the Serre--Tate expansion. The proof is sorry-backed at this stage.
\end{proof}

% SOURCE: [MR4372220], Thm.~\ref{thm:thmC}, \S{}Proof of the $p$-part of BSD formula (read from references/MR4372220-anticyclotomic-iwasawa-source/Eisenstein.tex)
% SOURCE QUOTE: "Let $E/\bQ$ be an elliptic curve, and let $p>2$ be a prime of good ordinary reduction for $E$. Assume that $E$ admits a cyclic $p$-isogeny with kernel $C=\mathbb{F}_p(\phi)$, with $\phi:G_\bQ\rightarrow\mathbb{F}_p^\times$ such that $\phi\vert_{G_p}\neq\mathds{1},\omega$, $\phi$ is either ramified at $p$ and odd, or unramified at $p$ and even. If ${\rm ord}_{s=1}L(E,s)=1$, then $\ord_{p}\bigl(\frac{L'(E, 1)}{\operatorname{Reg}(E / \Q) \cdot \Omega_{E}}\bigr)=\ord_{p}\bigl(\# \sha(E / \Q) \prod_{\ell \nmid \infty} c_{\ell}(E / \Q)\bigr)$. In other words, the $p$-part of the Birch--Swinnerton-Dyer formula for $E$ holds."
\begin{theorem}[$p$-part of the Birch--Swinnerton-Dyer formula]
  \label{thm:thmC}
  \lean{MR4372220.Applications.bsdPPart}
  \uses{thm:thmA, def:iwasawa-algebra, thm:control-jsw, thm:padic-gz}
  \textit{Source: MR4372220, \S{}Proof of the $p$-part of BSD formula, Theorem~C.}
  Let $E/\mathbb{Q}$ be an elliptic curve, and let $p > 2$ be a prime of good
  ordinary reduction for $E$. Assume that $E$ admits a cyclic $p$-isogeny with
  kernel $\mathbb{F}_p(\phi)$, where $\phi : G_{\mathbb{Q}} \to \mathbb{F}_p^\times$
  satisfies $\phi|_{G_p} \neq \mathbf{1}, \omega$, and $\phi$ is either ramified
  at~$p$ and odd, or unramified at~$p$ and even.
  If $\operatorname{ord}_{s=1} L(E, s) = 1$, then
  \[
    \operatorname{ord}_p\!\left(\frac{L'(E, 1)}{\operatorname{Reg}(E/\mathbb{Q}) \cdot \Omega_E}\right)
    =
    \operatorname{ord}_p\!\left(\#\sha(E/\mathbb{Q})\prod_{\ell \nmid \infty} c_\ell(E/\mathbb{Q})\right),
  \]
  where $\operatorname{Reg}(E/\mathbb{Q})$ is the regulator of $E(\mathbb{Q})$,
  $\Omega_E = \int_{E(\mathbb{R})} |\omega_E|$ is the N\'{e}ron period, and
  $c_\ell(E/\mathbb{Q})$ is the Tamagawa number of $E$ at~$\ell$.
  In other words, the $p$-part of the Birch--Swinnerton-Dyer formula for $E$ holds.
\end{theorem}

% SOURCE QUOTE PROOF: "Suppose ${\rm ord}_{s=1}L(E,s)=1$ and choose, as in the proof of Theorem~\ref{thm:thmB}, an imaginary quadratic field $K$ of discriminant $D_K$ such that (a) $D_K<-4$ is odd, (b) every prime $\ell$ dividing $N$ splits in $K$, (c) $p$ splits in $K$, say $p=v\bar{v}$, (d) $L(E^K,1)\neq 0$. Then ${\rm ord}_{s=1}L(E/K,s)=1$, which by Theorem~\ref{thmGZ} implies that the Heegner point $P_K\in E(K)$ has infinite order, and therefore ${\rm rank}_\Z E(K)=1$ and $\#\sha(E/K)<\infty$ by \cite{kolyvagin-mw-sha}. In particular, \eqref{eq:intro-sel} holds, and so all the hypotheses of Theorem \ref{thm:thmA} are satisfied. Thus there is a $p$-adic unit $u\in(\Z_p^{{\rm ur}})^\times$ for which $\mathcal{F}_E(0)=u\cdot\mathcal{L}_E(0)$, where $\mathcal{F}_E\in\Lambda$ is a generator of ${\rm char}_\Lambda(\mathfrak{X}_E)$. The hypotheses on $\phi$ imply that $E(K)[p]=0$, and so Theorem~\ref{controlthm} applies with $P=P_K$, which combined with Theorem~\ref{thmpadicGZ} [...] yields the equality ${\rm ord}_p\bigl(\#\sha(E/K)\bigr)=2\;{\rm ord}_p\bigl(c_E^{-1}u_K^{-1}\cdot[E(K):\Z.P_K]\bigr)-\sum_{w\in S}{\rm ord}_p\bigl(c_w(E/K)\bigr)$."
\begin{proof}
  \uses{thm:thmA, def:iwasawa-algebra}
  Assume $\operatorname{ord}_{s=1} L(E, s) = 1$. Choose an imaginary quadratic field $K$
  of odd discriminant $D_K < -4$ such that every prime $\ell \mid N$ splits in~$K$,
  $p = v\bar{v}$ splits in~$K$, and $L(E^K, 1) \neq 0$ (as in the proof of
  \cref{thm:thmB}, Case~$r = 1$).

  \textbf{Step 1: Apply \cref{thm:thmA}.}
  Then $\operatorname{ord}_{s=1} L(E/K, s) = 1$, so by the classical Gross--Zagier
  formula the Heegner point $P_K \in E(K)$ has infinite order, hence
  $\operatorname{rank}_{\mathbb{Z}} E(K) = 1$ and $\sha(E/K)[p^\infty]$ is finite
  by Kolyvagin. In particular, hypothesis~(Sel) holds, so all hypotheses of
  \cref{thm:thmA} are satisfied. There exists a $p$-adic unit
  $u \in (\mathbb{Z}_p^{\mathrm{ur}})^\times$ such that
  \[
    \mathcal{F}_E(0) = u \cdot \mathcal{L}_E(0),
  \]
  where $\mathcal{F}_E \in \Lambda$ generates $\mathrm{char}_\Lambda(\mathfrak{X}_E)$.

  \textbf{Step 2: Anticyclotomic control theorem.}
  The hypotheses on $\phi$ imply $E(K)[p] = 0$. The anticyclotomic control
  theorem (Theorem~3.3.1 of Jetchev--Skinner--Wan) applies with $P = P_K$:
  \[
    \#\mathbb{Z}_p/\mathcal{F}_E(0)
    = \#\sha(E/K)[p^\infty] \cdot
      \left(\frac{\#(\mathbb{Z}_p/((1{-}a_p{+}p)/p)\cdot\log_{\omega_E} P_K))}{[E(K):\mathbb{Z}{\cdot}P_K]_p}\right)^2
      \cdot \prod_{w \mid N} c_w(E/K)_p.
  \]

  \textbf{Step 3: $p$-adic Gross--Zagier formula.}
  By the BDP interpolation formula,
  $\mathcal{L}_E(0) = c_E^{-2}(1 - a_p p^{-1} + p^{-1})^2 \log_{\omega_E}(P_K)^2$.
  Combining Steps~1--3 gives
  \[
    \operatorname{ord}_p\!\bigl(\#\sha(E/K)\bigr)
    = 2\operatorname{ord}_p\!\bigl(c_E^{-1}[E(K):\mathbb{Z}{\cdot}P_K]\bigr)
      - \sum_{w \in S} \operatorname{ord}_p(c_w(E/K)).
  \]

  \textbf{Step 4: Classical Gross--Zagier formula and comparison.}
  The classical Gross--Zagier formula (rewritten as in \cite[p.~312]{grosszagier})
  and the factorization $L(E/K, s) = L(E, s) L(E^K, s)$ give
  \[
    \operatorname{ord}_p\!\left(\frac{L'(E,1)}{\operatorname{Reg}(E/\mathbb{Q})\cdot\Omega_E \cdot \prod_\ell c_\ell(E/\mathbb{Q})}\right)
    - \operatorname{ord}_p(\#\sha(E/\mathbb{Q}))
    = \operatorname{ord}_p\!\left(\frac{L(E^K,1)}{\Omega_{E^K} \cdot \prod_\ell c_\ell(E^K/\mathbb{Q})}\right)
    - \operatorname{ord}_p(\#\sha(E^K/\mathbb{Q})).
  \]
  The hypotheses on $\phi$ ensure that $E^K$ satisfies the conditions of
  Greenberg--Vatsal (the Euler factor character is ramified-even or unramified-odd),
  and since $L(E^K, 1) \neq 0$, the Greenberg--Vatsal theorem shows that the
  right-hand side vanishes, completing the proof.
\end{proof}
